% Source: https://artofproblemsolving.com/community/c3h1435941_number_theory_modular_stuff
% Stable identity: aops-topic-1435941-q01
From the second congruence, $x=5+8t$. Substitute into the first: $5+8t\equiv2\pmod7$. Since $8\equiv1\pmod7$, this is $5+t\equiv2$, so $t\equiv-3\equiv4\pmod7$, that is $t=4+7s$. Then
\[
x=5+8(4+7s)=37+56s .
\]
So the solutions are exactly the integers $x\equiv37\pmod{56}$, and the smallest positive one is $\boxed{37}$ (check: $37=5\cdot7+2=4\cdot8+5$).

The efficient method: write the solution of the congruence with the larger modulus as $x=5+8t$, reduce the other congruence to a congruence for $t$, and solve that one. Since $7$ and $8$ are coprime, the answer is unique modulo $7\cdot8=56$ (Chinese remainder theorem).
